\chapter{Generative AI Usage}
Claude (Anthropic, version Sonnet 4.5) was utilized as a writing assistant throughout this project for two primary purposes: enhancing the linguistic quality and academic tone of technical documentation, and formatting Alloy formal specification code for improved readability.

\section{Linguistic Enhancement}

\textbf{AI Tool Used:} Claude (Anthropic, Sonnet 4.5)

\textbf{Purpose:} Refining user scenario descriptions, polishing glossary definitions, and improving the clarity of system component explanations while maintaining technical precision.

\textbf{Typical Usage Pattern:} Initial drafts were written by the authors and then submitted to the AI with prompts requesting linguistic refinement and academic tone optimization. The AI provided alternative phrasings and structural suggestions.

\textbf{Example Prompt:}
\begin{verbatim}
"Optimize the following phrase in English. 
Target: brilliant academic style"
\end{verbatim}

\textbf{Before:} "The Registered User David manually inserts information about a specific bike path by inserting one street names and its status at a time (e.g., "optimal", "requires maintenance") without the need of real world biking and automatic recording."

\textbf{After:} "The registered user David manually inserts information about a specific bike path by entering individual street names and their status (e.g., "optimal", "requires maintenance") without requiring physical cycling or automated recording."

\textbf{Key Improvements:}
\begin{itemize}
    \item Grammatical corrections ("one street names" → "individual street names")
    \item Consistency in capitalization ("Registered User" → "registered user")
    \item Conciseness ("without the need of" → "without requiring")
    \item Professional terminology ("real world biking" → "physical cycling")
\end{itemize}

\textbf{Human Verification:} All AI-generated outputs were carefully reviewed, modified, and validated by the authors to ensure correctness and consistency with the BBP system specifications.

\section{Alloy Code Formatting}

\textbf{AI Tool Used:} Claude Code (Claude Sonnet 4.5)

\textbf{Purpose:} Formatting Alloy formal specification code to enhance readability, as no automated code formatters for Alloy are available in VSCode.

\textbf{Prompt Given:}
\begin{verbatim}
"Format the following alloy document. Do not change content, 
just tidy the elements placing"
\end{verbatim}

\textbf{AI Actions Performed:}
\begin{enumerate}
    \item \textbf{Syntax Error Correction:} Fixed missing closing braces in \texttt{PathLifecycle} and \texttt{ObstacleStatus} enum definitions
    \item \textbf{Consistent Indentation:} Standardized indentation throughout the document (2-space indentation, replacing mixed tabs/spaces)
    \item \textbf{Code Organization:} Added clear section dividers with descriptive headers:
    \begin{itemize}
        \item ENUMERATIONS
        \item SIGNATURES
        \item PUBLICATION FACTS AND FUNCTIONS
        \item USER AND TRIP OWNERSHIP FACTS
        \item OBSTACLE LIFECYCLE
        \item MANUAL TRIP LIFECYCLE
        \item AUTOMATIC TRIP LIFECYCLE
        \item PATH LIFECYCLE + STATUS EVOLUTION
    \end{itemize}
    \item \textbf{Alignment:} Aligned field declarations, inline comments, and code blocks consistently
    \item \textbf{Whitespace Normalization:} Removed inconsistent blank lines and standardized spacing
\end{enumerate}

\textbf{Outcome:} The Alloy specification became more readable and maintainable while preserving all original logic, semantics, and functionality. The improved structure facilitated easier navigation and understanding of the formal model's different components.

\textbf{Human Verification:} All formatting changes were reviewed and the formatted code was validated to ensure no semantic changes were introduced.

\vspace{1em}
\noindent\textit{Note:} All substantive technical content, system design decisions, architectural choices, and UML diagrams remain the original work of the author(s). AI assistance was limited to linguistic refinement and code formatting only.